Many important physical theories are described by Lagrangians that are invariant under local symmetry transformations that form Lie groups; such theories are named {\it gauge theories}.
For example, the Standard Model of particle physics, which is our most accurate description of Nature at the shortest length-scales, is a quantum field theory centered around the action of three gauge groups~\cite{Weinberg:1967,Salam:1968rm,Greenberg:1964,Han:1965}, and several important condensed matter systems can be described by effective gauge theories~\cite{Kogut1979SpinSystems, Seiberg2016GaugeCMTDuality,Sachdev2018EmergentGauge,Guo2020EmergentGauge}.
In the strong-coupling limit, these theories are non-perturbative, and numerical formulations on discrete spacetime lattices offer the only known way to compute properties of interest from first principles.

Calculations within lattice frameworks typically proceed by estimating expectation values of observables using Markov Chain Monte Carlo (MCMC) to sample from thermodynamic distributions or Euclidean-time path integrals. In both cases, samples $U$ (typically high-dimensional) are drawn from an exponentially weighted distribution $p(U) = e^{-S(U)} / Z$, where the physics is encoded in an energy or action functional $S(U)$, and the normalizing constant $Z$ is unknown. When MCMC sampling from the distribution $p(U)$ is efficient, precise physical predictions can be made from the theory. However, as the model parameters are tuned towards criticality, e.g.~to describe universal properties of condensed matter theories or to access the continuum limit of quantum field theories, critical slowing down (CSD) can cause the computational cost of sampling to diverge~\cite{Wolff:1989wq}.

Specialized approaches have been developed to avoid CSD for specific theories~\cite{Kandel:1988zz,Bonati:2017woi,Wolff:1988uh,Hasenbusch:2017unr,Hasenbusch:2018ztj,Swendsen:1987ce,Prokofev:2001ddj,Kawashima2004,Bietenholz:1995zk,Luscher:2011kk}. For several theories of interest, however, CSD obstructs calculations. This is true in particular for the lattice formulation of quantum chromodynamics (QCD)~\cite{Aoki:2006,Schaefer:2009xx,McGlynn:2014}, which enables calculations of non-perturbative phenomena arising from the Standard Model of particle physics. Recently, there has been progress in the development of \emph{flow-based} generative models which can be trained to directly produce samples from a given probability distribution; early success has been demonstrated in theories of bosonic matter, spin systems, molecular systems, and for Brownian motion~\cite{LiWang2018NNRG,ZhangEWang2018Monge,Albergo2019FlowMCMC,Hartnett:2020,khler2019equivariant,Noeeaaw1147,li2019neural,madhawa2019graphnvp,liu2019graph,shi2020graphaf,both2019temporal}. This progress builds on the great success of flow-based approaches for image, text, and structured object generation~\cite{Dinh:2014,Dinh:2016,Kingma:2018,Ho:2019, tran2019discrete,ziegler2019latent,kumar2019videoflow,gu2019pointflow}, as well as non-flow-based machine learning techniques applied to sampling for physics~\cite{Huang:2017,Wang2017,Wu:2019,zhang2019chem, Carrasquilla_2019}. If flow-based algorithms can be designed and implemented at the scale of state-of-the-art calculations, they would enable efficient sampling in lattice theories that are currently hindered by CSD.

\begin{figure}[t]
\includegraphics{images/sample_hmc_topo_history.pdf}
\caption{Standard approaches (HMC and HB) to MCMC sampling for $\Uone$ gauge theory explore the distribution of topological charge $Q$ very slowly compared with the flow-based approach introduced here. Results are shown for coupling $\beta=7$ on a $16\times 16$ lattice, see Eq.~\eqref{eq:action}. The first (second) half of the Markov chain history is displayed for HMC (HB).}
\label{fig:topo_freezing}
\end{figure}

In this Letter, we develop a provably correct flow-based sampling algorithm designed for lattice gauge theories, including lattice QCD. We demonstrate the application of this approach to $\Uone$ gauge theory in two spacetime dimensions. This theory is solvable, and thus provides a testing ground where the accuracy of numerical methods can be checked. Two standard MCMC approaches, Hybrid Monte Carlo (HMC)~\cite{Duane1987HMC} and Heat Bath (HB)~\cite{Creutz:1979dw,Cabibbo:1982zn,Kennedy:1985nu}, suffer from critical slowing down in this theory; for example, Fig.~\ref{fig:topo_freezing} depicts Markov chain histories for sampling near the continuum limit, in which both methods explore topological sectors very slowly. Using our flow-based algorithm, independent samples of field configurations are produced with appropriate frequency from each topological sector, enabling far more accurate estimation of topological quantities at a given computational cost. Critical to the success of this approach is enforcing exact gauge symmetry in the flow-based distribution: when the symmetry is enforced, we can successfully train flow-based models at a range of parameters approaching criticality, while without it, models of similar scale fail to learn the distributions under the same training procedure.
